\documentclass{article}
\usepackage[utf8]{inputenc}
\usepackage[spanish]{babel}
\usepackage[letterpaper,top=2cm,bottom=2cm,left=3cm,right=3cm,marginparwidth=1.75cm]{geometry}
\usepackage{float}

\usepackage{amsmath}
\usepackage{amssymb}
\usepackage{graphicx}


\title{Tarea 07 de Lenguajes de Programación}
\author{
  Ramírez Ríos, Víctor Hugo\\
  \texttt{a221205528@unison.mx}
}
\date{Noviembre 2023}


\begin{document}

\maketitle

\begin{enumerate}
    \item Extiende el lenguaje agregando un nuevo operador minus que toma un argumento $n$ y regresa $-n$. Por ejemplo, el valor de minus(-(minus(5),9)) debe ser 14.
    \begin{itemize}
        \item Léxico \\
        minus
        \item Sintaxis 
        \begin{itemize}
            \item Concreta \\
            $Expression$$\to$minus($Expression$)
            \item Abstracta \\
            (minus-exp $exp1$)
        \end{itemize}
        \item Semántica \\
        (value-of(minus-exp $exp1$)$\rho$) = (num-val (- (expval$\to$num (value-of $exp1$ $\rho$))
    \end{itemize}
    \item Extiende el lenguaje agregando operadores para la suma, multiplicación y cociente de enteros.
    \begin{itemize}
        \item Léxico \\
        + \\
        * \\
        / 
        \item Sintaxis 
        \begin{itemize}
            \item Concreta \\
            $Expression$$\to$+($Expression$, $Expression$) \\
            $Expression$$\to$*($Expression$, $Expression$) \\
            $Expression$$\to$/($Expression$, $Expression$) 
            \item Abstracta \\
            (sum-exp $exp1$ $exp2$) \\
            (mult-exp $exp1$ $exp2$) \\
            (div-exp $exp1$ $exp2$)
        \end{itemize}
        \item Semántica 
        \begin{align*}
            (\text{value-of}(\text{sum-exp}& exp1 exp2) \rho ) \\
            =(\text{num-val} (+ &(\text{expval} \to \text{num} (\text{value-of} exp1 \rho))\\
            &(\text{expval} \to \text{num} ( \text{value-of} exp1 \rho))))
        \end{align*}
        
        \begin{align*}
            (\text{value-of}(\text{mult-exp}& exp1 exp2) \rho ) \\
            =(\text{num-val} (* &(\text{expval} \to \text{num} (\text{value-of} exp1 \rho))\\
            &(\text{expval} \to \text{num} ( \text{value-of} exp1 \rho))))
        \end{align*}
        
        \begin{align*}
            (\text{value-of}(\text{div-exp}& exp1 exp2) \rho ) \\
            =(\text{num-val} (\text{floor}(/ &(\text{expval} \to \text{num} (\text{value-of} exp1 \rho))\\
            &(\text{expval} \to \text{num} ( \text{value-of} exp1 \rho)))))
        \end{align*}
    \end{itemize}
    \item Agrega un predicado de igualdad numérica equal? y predicados de orden greater? y less? al conjunto de operaciones del lenguaje LET.

    \begin{itemize}
        \item Léxico \\
        equal? \\
        greater? \\
        less? 
        \item Sintaxis 
        \begin{itemize}
            \item Concreta \\
            $Expression$$\to$equal?($Expression$, $Expression$) \\
            $Expression$$\to$greater?($Expression$, $Expression$) \\
            $Expression$$\to$less?($Expression$, $Expression$) 
            \item Abstracta \\
            (equal?-exp $exp1$ $exp2$) \\
            (greater?-exp $exp1$ $exp2$) \\
            (less?-exp $exp1$ $exp2$)
        \end{itemize}
        \item Semántica 
        \begin{align*}
            (\text{value-of}(\text{equal?-exp}& exp1 exp2) \rho ) \\
            =(\text{num-val} (equal? &(\text{expval} \to \text{num} (\text{value-of} exp1 \rho))\\
            &(\text{expval} \to \text{num} ( \text{value-of} exp1 \rho))))
        \end{align*}
        
        \begin{align*}
            (\text{value-of}(\text{greater?-exp}& exp1 exp2) \rho ) \\
            =(\text{num-val} (>  (&\text{expval} \to \text{num} (\text{value-of} exp1 \rho))\\
            &(\text{expval} \to \text{num} ( \text{value-of} exp1 \rho))))
        \end{align*}
        
        \begin{align*}
            (\text{value-of}(\text{less?-exp}& exp1 exp2) \rho ) \\
            =(\text{num-val} (< (&\text{expval} \to \text{num} (\text{value-of} exp1 \rho))\\
            &(\text{expval} \to \text{num} ( \text{value-of} exp1 \rho))))
        \end{align*}
    \end{itemize}
    \item Agrega operaciones de procesamiento de listas al lenguaje, incluyendo cons, car, cdr, null? y emptylist.
    
    \begin{itemize}
        \item Léxico \\
        cons \\
        car \\
        cdr \\ 
        null? \\ 
        emptylist  
        \item Sintaxis 
        \begin{itemize}
            \item Concreta \\
            $Expression$$\to$cons($Expression$, $Expression$) \\
            $Expression$$\to$car($Expression$) \\
            $Expression$$\to$cdr($Expression$) \\
            $Expression$$\to$null?($Expression$) \\
            $Expression$$\to$emptylist() 
            \item Abstracta \\
            (cons-exp $exp1$ $exp2$) \\
            (car-exp $exp1$) \\
            (cdr-exp $exp1$) \\
            (null?-exp $exp1$) \\
            (emptylist-exp $exp1$)
        \end{itemize}
        \item Semántica 
        \begin{align*}
            (\text{value-of}(\text{cons-exp}& exp1 exp2) \rho ) \\
            =(\textbf{let} (&[val1 (\text{value-of} exp1 \rho)] \\
                         &[val2 (\text{value-of} exp2 \rho)]) \\
                         & (\text{pair-val} (\text{cons} val1 val2)))
        \end{align*}
        
        \begin{align*}
            (\text{value-of}(\text{car-exp}& exp1) \rho ) \\
            =(\text{pair-val} (\text{car} (&\text{expval} \to \text{num} (\text{value-of} exp1 \rho))\\
        \end{align*}
        
        \begin{align*}
            (\text{value-of}(\text{cdr-exp}& exp1) \rho ) \\
            =(\text{pair-val} (\text{cdr} (&\text{expval} \to \text{num} (\text{value-of} exp1 \rho))\\
        \end{align*}

        \begin{align*}
            (\text{value-of}(\text{null?-exp} \; & exp1) \; \rho ) \\
            =(\textbf{let} (&[val1 (\text{value-of} \; exp1 \; \rho)]) \\
                         & (\text{bool-val} (= \text{null-val} \; val1)))
        \end{align*}

        \begin{align*}
            (\text{value-of}(\text{emptylist}& exp1) \rho ) \\
            =(\text{null-val})\\
        \end{align*}
    \end{itemize}
    \item Agrega una operacion list al lenguaje. Esta operacion debe tomar cualquier cantidad de argumentos y regresar un valor expresado de la lista de sus valores.
    
    \begin{itemize}
        \item Léxico \\
        list
        \item Sintaxis 
        \begin{itemize}
            \item Concreta \\
            $Expression$ $\to$ list($Expressions$)  \\
            $Expressions$ $\to$ $\epsilon$ \\
            $Expressions$ $\to$ $ExpressionsAux$ \\
            $ExpressionsAux$ $\to$ $Expression$ \\
            $ExpressionsAux$ $\to$ $Expression, ExpressionsAux$
            \item Abstracta \\
            (list-exp $exps$) 
        \end{itemize}
        \item Semántica 
        \begin{align*}
            (\text{value-of}(\text{list-exp} \; & exps) \rho ) \\
            =(\textbf{let} (&[val1 (\text{value-of} \; exps \; \rho)]) \\
                          (\text{if} &(\text{expval} \to \text{bool} (= \text{null-val} \; val1)) \\
                         &(\text{null-val)}) \\
                         &(\text{haz una lista con el cons, car y cdr}))) \\
        \end{align*}
    \end{itemize}
    \item Incorpora al lenguaje expresiones cond. Usa la gramatica
    \[ Expression \to \textbf{cond}\{ Expression => Expression\}* \textbf{end}\]

    \begin{itemize}
        \item Léxico \\
        cond \\
        end \\
        =$>$
        \item Sintaxis 
        \begin{itemize}
            \item Concreta \\
                    \[ Expression \to \textbf{cond}\{ Expression => Expression\}* \textbf{end}\]
            \item Abstracta \\
            (cond-exp $exps$) 
        \end{itemize}
        \item Semántica 
    \end{itemize}
    \item Cambia los valores del lenguaje para que los enteros sean los unicos valores expresados. Modifica if para que el valor de 0 sea tratado como falso y todos los otros valores sean tratados como verdaderos. Modifica los predicados de manera consistente.

        \begin{align*}
        \varepsilon(\textbf{zero?}(exp1),\rho)= \begin{cases}
       1 &\quad\text{si} \; \varepsilon(exp1, \rho) = 0\\
       0 &\quad\text{si} \; \varepsilon(exp1, \rho) \neq 0\\ 
     \end{cases}
        \end{align*}

        \begin{align*}
        \varepsilon(\textbf{if}(exp1)\textbf{then}(exp2)\textbf{else}(exp3),\rho)= \begin{cases}
       \varepsilon(exp2, \rho) &\quad\text{si} \; \varepsilon(exp1, \rho) = 1\\
       \varepsilon(exp3, \rho) &\quad\text{si} \; \varepsilon(exp1, \rho) = 0 \\
     \end{cases}
        \end{align*}
    
    
    \item Como una alternativa al ejercicio anterior, agrega una nueva categoria sintactica Bool-exp de expresiones booleanas al lenguaje. Cambia la producción para expresiones condicionales para que sea
    \[ Expression \to \textbf{if} \; \text{Bool-exp} \; \textbf{then} expression \textbf{else} Expression \]

    Escribe producciones apropiadas para Bool-exp y especifica su semantica con value-of-bool-exp. En donde terminan estando los predicados del ejercicio 3 con este cambio?

    \begin{itemize}
        \item Léxico \\
        \item Sintaxis 
        \begin{itemize}
            \item Concreta \\
        \[ Expression \to \textbf{if} \; \text{Bool-exp} \; \textbf{then} expression \textbf{else} Expression \]
        \[\text{Bool-exp} = 1\]
        \[\text{Bool-exp} = 0\]
        \[\text{Bool-exp} = zero?(Expression)\]
        \[\text{Bool-exp} = equal?(Expression, Expression)\]
        \[\text{Bool-exp} = greater?(Expression, Expression)\]
        \[\text{Bool-exp} = 0\]
            \item Abstracta \\
            (if boolexp exp1 exp2) \\
            (bool-exp $bool$)  \\
            (zero?-boolexp $exp1$)  \\
            (equal?-boolexp $exp1$ $exp2$)  \\
            (greater?-boolexp $exp1$ $exp2$) \\
            (less?-boolexp $exp1$ $exp2$) \\
        \end{itemize}
        \item Semántica 
        
        \begin{align*}
            (\text{value-of}(\text{bool-exp}& \; bool) \rho ) = (\text{bool-val} \; bool) \\
        \end{align*}


        \begin{align*}
        \varepsilon(\textbf{if}(boolexp)\textbf{then}(exp1)\textbf{else}(exp2),\rho)= \begin{cases}
       \varepsilon(exp1, \rho) &\quad\text{si} \; \varepsilon(boolexp, \rho) = T\\
       \varepsilon(exp2, \rho) &\quad\text{si} \; \varepsilon(boolexp, \rho) = F \\
     \end{cases}
        \end{align*}
    \end{itemize}

    \item Modifica la implementacion del interprete agregando una nueva operacion print que toma un argumento, lo imprime y regresa el entero 1. Por que esta expresion no es expresable en nuestro metodo de especificacion formal?

    Por que no tenemos una forma de representar la operacion de impresión usando la especificación formal.

    \item Extiende el lenguaje para que las expresiones let puedan vincular una cantidad arbitraria de variables, usando la producción, 
    \[Expression \to \textbf{let} \{Identifier = Expression\}* \textbf{in} \; Expression\]

    \begin{itemize}
        \item Sintactica
        \begin{itemize}
            \item Concreta
            
    \[Expression \to \textbf{let} \{Identifier = Expression\}* \textbf{in} \; Expression\]
            \item Abstracta \\
            (let-exp (var exp)* body)
        \end{itemize}
        \item Semantica
        \begin{align*}
            \frac{\varepsilon(exp1, \rho) = val1 \: \: \: \varepsilon(exp2, \rho) = val2 \dots}{\varepsilon(\textbf{let} \; x1 = exp1 \: \: \: x2 = exp2 \dots \textbf{in} \; body, \rho) = \varepsilon (body, [x1 : val1, x2 : val2 \dots] \rho)}
        \end{align*}
    \end{itemize}

    
    
    \item Extiende el lenguaje con una expresión let* que funciona como en Racket.

    

    \item Agrega una expresión al lenguaje de acuerdo a la siguiente regla, 
    \[Expression \to \textbf{unpack} \{Identifier\}*  = Expression \textbf{in} Expression\]

    \item Consideremos funciones de compilación compile-program y compile que toman de argumento un árbol de sintaxis y regresan un árbol de sintaxis. Implementa estas funciones para que el lenguaje compile el cond del ejercicio 7, la extensión a let del ejercicio 11, y el let* del ejercicio 12, utilizando el lenguaje LET base y las extensiones del lenguaje 4,5 y 13.

    
\end{enumerate}
\end{document}